\section{Discussion}
\label{sec:discussion}

In scientific machine learning, $\Delta$-learning rests on a compelling intuition: rather than predicting a high-fidelity quantum-mechanical energy from scratch, a GNN should learn the correction to a cheap approximate baseline. Because the baseline captures much of the variation in the target, the resulting residual is smaller and presumed easier to learn. Yet, as our findings demonstrate, this intuition does not always hold. Aggressively shrinking total target scale does not guarantee downstream learnability, explaining why $\Delta$-learning has often relied on trial-and-error heuristics rather than design principles.

Our results resolve this inconsistency by reframing $\Delta$-learning as a functional target transformation governed by two distinct properties: \textbf{target scale} and \textbf{scale-normalized representation smoothness}, which tracks downstream learnability as assessed by $n\text{-MAE}$. Because absolute test error is the product of target scale and normalized error, aggressive scale reduction can still yield a net accuracy gain if the shrinkage outweighs the accompanying loss in learnability—as seen with the linear Bag-of-Elements (BoE) baseline. However, pushing shrinkage even further with a nonlinear MLP-BoE baseline raises normalized error enough to outweigh the further reduction in target scale and degrades performance.

Optimal $\Delta$-learning therefore requires \textbf{baseline complementarity}: we propose prioritizing baselines that capture the specific components responsible for the model's error, rather than merely absorbing the largest share of total target spread. Local descriptor baselines (Bag-of-Bonds, PRE, and PRE+A) share geometric information with local message-passing GNNs; they strip away smooth local features the network already represents efficiently, leaving residuals that are narrower in scale yet disproportionately rough and harder to learn, resulting in poorer test performance than direct learning. Conversely, semi-empirical quantum methods (xTB and DFTB) capture electronic effects—such as charge redistribution and long-range electrostatics—that may contribute to the errors of local architectures. By absorbing the specific energy contributions that drive GNN error, complementary baselines shrink target scale while preserving or reducing scale-normalized roughness, improving test accuracy by 30–40\% and substantially enhancing energetic rank correlation across distinct chemical domains (QM9).

Because post-training evaluation requires expensive model optimization, practical target engineering requires pre-training diagnostics. To address this, we introduced the scale-normalized graph Dirichlet roughness ($D_{\rm IQR}$), which tracks target variation over architecture-aligned representation proxies without training a downstream GNN. Its strong semi-log association with normalized test error ($n\text{-MAE}$) supports $D_{\rm IQR}$ as a candidate pre-training diagnostic for screening baselines. While shared normalization contributes substantially to this association, joint-label permutation controls support an additional contribution from molecular-neighbourhood alignment.

\section{Conclusion}
\label{sec:conclusion}

This study advances $\Delta$-learning from heuristic baseline selection toward principled target design. By establishing that baseline complementarity and the balance between target shrinkage and residual smoothness govern learnability, our results position target formulation as a foundational design choice alongside network architecture, offering a clear path forward for residual learning, multi-fidelity modeling, and scientific machine learning surrogates.
